\subsection{Ordinary surfaces with arbitrary polarization degree}
\label{sec:DE-A3-ordinary-all-polarizations}

\begin{lemma}[The ordinary polarization kernel and its connected isotropic part]
\label{lem:DE-A3-ordinary-polarization-kernel}
\status{Proved; root and reduction-descent full independent proof/source audits passed.}
Let $A$ be an ordinary abelian surface over an algebraically closed
field of characteristic two, and let $M$ be ample with $D=\chi(M)$.
Then $\ker\lambda_M$ is killed by $D$.  If $D_2$ is the largest
power of two dividing $D$, the two-primary part of this kernel is
\[
             K_2=K_2^0\times K_2^{\et},
             \qquad \rank K_2^0=\rank K_2^{\et}=D_2,
\]
where $K_2^0$ is multiplicative and $K_2^{\et}$ is constant.
The commutator pairing is trivial on either factor and pairs the
factors perfectly.  In particular $K_2^0$ is maximal isotropic
in $K_2$.
\end{lemma}
\begin{proof}
We first justify the group-scheme assertions without a separability
assumption.  Over a perfect field, a finite commutative group scheme
splits into its connected and etale parts; there are no nonzero
homomorphisms between a finite local and an etale group scheme.
These are \cite[Lemma~4.44 and Proposition~4.45, p.~67]
{EdixhovenVanDerGeerMoonenAV}, including their proofs.
Ordinarity, the prime-power torsion count, and Cartier duality give
\[
              A[2^n]\simeq\mu_{2^n}^{2}
                      \times(\mathbf Z/2^n\mathbf Z)^2.
\]
For completeness, the etale part has order $2^{2n}$ by
\cite[Proposition~5.22, pp.~80--81]{EdixhovenVanDerGeerMoonenAV}.
The dual surface is isogenous to $A$, hence is ordinary by that
same proposition.  The Cartier dual of its etale torsion is a
multiplicative subgroup of $A[2^n]$ of order $2^{2n}$, by
\cite[Corollary~7.17, p.~104]{EdixhovenVanDerGeerMoonenAV}.
This is the entire connected part, since $\deg[2^n]=2^{4n}$.
Thus every finite two-primary subgroup of $A$ has multiplicative
connected part and constant etale part.

Put $K=\ker\lambda_M$.  Its commutator pairing is perfect and
alternating, and $\rank K=D^2$; see
\cite[Proposition~8.19 and Corollary~8.20, p.~119]
{EdixhovenVanDerGeerMoonenAV}.
Distinct primary components are orthogonal.  On $K_2$, the pairing
on the connected factor is trivial, because it gives a homomorphism
from a multiplicative local group to its etale Cartier dual.
The pairing on the etale factor is likewise trivial.  Perfection
therefore identifies each factor with the Cartier dual of the other.
Their orders are equal and have product $D_2^2$, proving the claimed
orders and the maximal isotropy.  Both factors are killed by $D_2$:
this follows from the elementary exponent bound for a finite abelian
group of order $D_2$, applied also to the character group of the
multiplicative factor.

For an odd prime $\ell$, the $\ell$-primary part of $K$ is a
constant finite abelian group with a perfect alternating pairing.
If $x$ has maximal order $\ell^b$, perfection supplies $y$ such
that $e_M(x,y)$ has that same order.  Pairing a relation
$ux+vy=0$ successively with $x,y$ shows that $\langle x,y\rangle$
has order $\ell^{2b}$.  Consequently
$2b\le v_\ell(\rank K)=2v_\ell(D)$.
The exponent of this primary part thus divides the $\ell$-part
of $D$.  Combining all primary parts proves $DK=0$.
\end{proof}

\begin{lemma}[An etale principal norm source in the ordinary case]
\label{lem:DE-A3-ordinary-etale-norm}
\status{Proved; root and reduction-descent full independent proof/source audits passed.}
With $A,M,D$ as above, there is an ordinary principally polarized
surface $(B,\Theta_B)$ and a finite etale isogeny $\phi:B\to A$
of degree $D$ satisfying
\[
                 \phi^*M\equiv D\Theta_B,
                 \qquad \operatorname{Nm}_\phi(\Theta_B)\equiv M.
\]
The integer $D$ may be even.
\end{lemma}
\begin{proof}
The preceding lemma makes the restriction of $e_{M^D}=e_M^D$
to $\ker\lambda_M$ trivial.  The actual group-scheme descent
criterion \cite[Corollary~8.11 and proof, p.~117]
{EdixhovenVanDerGeerMoonenAV} therefore gives an ample line bundle
$N$ on $A^\vee$ such that $\lambda_M^*N\simeq M^D$.
This criterion allows nonreduced finite isotropic subgroup schemes.
As in Lemma~\ref{lem:DE-A3-principal-norm-source}, the pullback
polarization and degree identities give
\[
                       \lambda_N=D\lambda_M^{-1},
                       \qquad \chi(N)=D.
\]
The left side is an integral polarization, constructed by descent.

In the two-primary part of $\ker\lambda_N$, choose its entire
connected factor; it is maximal isotropic by
Lemma~\ref{lem:DE-A3-ordinary-polarization-kernel} applied to
$A^\vee,N$.  In the odd-primary part choose a maximal isotropic
subgroup.  Their product $H$ is maximal isotropic of order $D$.
Its connected part is multiplicative, and its remaining part has
odd order.  Therefore its Cartier dual $H^D$ is etale.
Proposition~8.19 of the same source gives a principal quotient
\[
               \psi:A^\vee\longrightarrow C=A^\vee/H,
               \qquad \psi^*\Theta_C\simeq N.
\]
Take $B=C^\vee$ with its dual principal polarization, and put
$\phi=\psi^\vee$.  Isogeny duality identifies
$\ker\phi=H^D$ \cite[Theorem~7.5, p.~100]
{EdixhovenVanDerGeerMoonenAV}; hence $\phi$ is etale even when
$D$ is even.  The typed polarization identity is
\[
                \psi\lambda_M\phi=D\lambda_C^{-1}.
\]
It gives $\phi^*M\equiv D\Theta_B$.  Applying norm and cancelling
$D$ in the torsion-free Neron--Severi group gives the second equality.
Ordinarity is preserved under isogeny, so $B$ is ordinary.
\end{proof}

\begin{theorem}[All degrees for every polarization on an ordinary surface]
\label{thm:DE-A3-ordinary-all-ranks}
\status{Proved; root and reduction-descent full independent proof/source audits passed.}
Let $A$ be an ordinary abelian surface over a perfect field of
characteristic two.  Let $M$ be any ample line bundle, put
$D=\chi(M)$ and $L=M^3$.  Then $L$ is very ample and, for every
$d\ge1$, the geometric generic smooth member of $|L^d|$ is ordinary,
with genus $9Dd^2+1$.  In particular its relative stable rank equals
its relative coherent dimension:
\[
                              r_d=9Dd^2-1.
\]
For ambient hypersurfaces in the embedding given by $L$, the same
formula holds whenever
\[
               d\ge\max\{1,\operatorname{reg}(\mathcal I_A)-1\}.
\]
There is no restriction on the parity of $D$ or $d$.
\end{theorem}
\begin{proof}
Work over an algebraic closure.  Apply
Lemma~\ref{lem:DE-A3-ordinary-etale-norm} to $M$.
The cubic principal bundle $L_B=\OO_B(3\Theta_B)$ is very ample.
Its norm has numerical class $M^3$ and is very ample by the
point-and-tangent norm-section argument in the proof of
Theorem~\ref{thm:DE-A3-separable-ranks}.  That argument requires
only a finite etale morphism, not that its degree be invertible.
Indeed one selects a section vanishing at one chosen lift, with
nonzero derivative when needed, and avoids the finitely many value
hyperplanes at every other lift.  The remaining factors of the norm
are units.  All ample line bundles in the same numerical class on
an abelian variety are translates, by surjectivity of the polarization.
Thus $L=M^3$ is very ample.

Theorem~\ref{thm:DE-A3-ppav-ranks} for the ordinary principal source
provides a nonempty open of smooth ordinary genus-ten curves
$C\in|L_B|$.  A general such curve has integral nodal image under
$\phi$.  We spell out why even kernel order creates no exception
to the nodal-image proof already used in
Lemma~\ref{lem:DE-A3-ppav-odd-norm}.
The kernel is a finite constant group acting freely by translations.
A nonzero element that preserves $L_B$ has a nontrivial projective
action in its very ample embedding, hence has a proper locus of
invariant sections; no semisimplicity or averaging is used.
A translation not preserving $L_B$ cannot preserve a member.
Thus a general $C$ has trivial kernel stabilizer.

For an indecomposable principal source, the two-first-jet failure
locus for a fixed nonzero kernel point $g$ is contained in
\[
                   3P+g\in\Theta_B,
                   \qquad 3P+2g\in\Theta_B.
\]
These are inverse images under $[3]$ of two distinct theta translates.
Their intersection is finite, including when $g$ has order two.
Away from these finitely many points, the two vanishing values and
proportionality of two nonzero gradients impose three independent
conditions while $P$ varies in dimension two.  At an exceptional
point, the two value conditions already define a proper locus.
The three-value separation lemma for $|3\Theta_B|$ similarly
excludes triple identifications for every pair of distinct kernel
points.  For a decomposable principal source, the cubic Segre
argument in the cited nodal-image proof supplies the same conclusion:
the one-coordinate coincidence has rank-five first jets, and the
three-collinear-point exception has dimension one and value rank two.
Those incidence bounds also depend only on distinctness of points,
not on the order of a translation.  Hence a general chosen ordinary
$C$ has only transverse double identifications and no triples.

Its image $Z$ is integral nodal, with normalization $C$.  Counting
the ordered pairs gives $18(D-1)$, so $Z$ has $9(D-1)$ nodes.
This count is a count of distinct geometric points under a free
etale action; division by two here counts unordered pairs and does
not invert two in the ground field.  The normalization sequence
adds a Frobenius-bijective graph space.  Consequently $Z$ is
ordinary, of arithmetic genus $10+9(D-1)=9D+1$.

The line bundle of $Z$ is $\operatorname{Nm}_\phi L_B$, whose
quotient with $M^3$ is in $\Pic^0(A)$.  Surjectivity of
$\lambda_{M^3}:A\to A^\vee$ supplies a single target translation
placing $Z$ in the exact system $|M^3|$.  This is an arbitrary
geometric translation; no inverse on the two-torsion points is
asserted.  Translation preserves ordinarity and the node structure.

Ample vanishing gives constant coherent dimensions for every divisor
in $|L^d|$, namely $9Dd^2+1$ in degree one on the divisor, and
$9Dd^2-1$ in the relative quotient.  Frobenius and its iterates
are maps of the corresponding bundles on the coefficient space.
Invertibility is open.  The ordinary nodal member just constructed
therefore yields a nonempty ordinary open for $d=1$, which meets
the nonempty smooth-member open.
For general $d$, choose $d$ smooth ordinary members of $|L|$
meeting pairwise transversely with no triple points.  Bertini and
avoidance of finitely many points give such a choice, and their
connected nodal union lies in $|L^d|$.  Its normalization and graph
terms are ordinary, so the union is ordinary.  The same open-rank
argument gives a smooth ordinary generic member of $|L^d|$.
Its ambient stable rank is two, and subtracting that ambient part
gives the displayed relative rank.  The final regularity statement
follows from surjectivity of ambient restriction in the stated range,
whose parameter fibres are integral, as in
Theorem~\ref{thm:DE-A3-separable-ranks}.
\end{proof}

\begin{remark}[Scope of the extension]
\label{rem:DE-A3-ordinary-all-rank-scope}
\status{Proved scope and verification statement.}
This theorem removes the parity restriction on the polarization degree
for ordinary surfaces.  Theorem~\ref{thm:DE-A3-ordinary-arbitrary-embedding}
further gives the exact eventual rank in any chosen ample embedding,
without requiring its line bundle to be a cube.  It does not assert a formula for arbitrary
nonordinary even-degree polarizations, or for arbitrary ambient varieties.
The finite group-scheme argument is essential: using an arbitrary
maximal isotropic subgroup could give a nonetale dual isogeny.

The retained program
\path{agents/obstructions_crystals/code/A3_ordinary_all_rank_checks.py}
checks the integral dual-polarization matrices, connected and etale
order identities, norm-node counts and normalization-graph arithmetic
for $211$ polarization types, including $146$ of even Euler
characteristic, and $27008$ degree cases.  Its author run passed in
$0.016162138999789022$ seconds.  The complete output is retained
alongside the program.  Root independently read the complete program
and reran it in $0.014947520998248365$ seconds, preserving the author
output and retaining its result at
\path{outputs/root_A3_ordinary_all_rank_checks.json}.
This finite arithmetic certificate does not
replace the group-scheme descent or geometric existence proofs.
\end{remark}
